\chapter{Factor Markets}\label{ch:b1:factor-markets}

In a town with one large employer, a higher wage floor can raise
employment. In a competitive labour market it cannot: a floor above the
market wage moves employers up their demand curve and they hire fewer
people. Which case holds depends on who sets the wage. This chapter
studies the markets for labour and capital from the buyer's side: what a
firm will pay for a factor, what happens when there is one buyer, and what
sets the share of income that goes to each factor.

\begin{openingchart}{Labour's share of GDP against GDP per person at
purchasing-power parity, 140 countries, 2019.}{Penn World Table 11.0
(labsh, rgdpe, pop)}
\begin{tikzpicture}
\begin{axis}[omscatter, xmode=log, log ticks with fixed point, xtickten={3,4,5}, xmin=200, xmax=200000, ymin=0, ymax=0.85,
  xlabel={GDP per person, PPP dollars (log scale)}, ylabel={Labour share of GDP},
  width=0.74\linewidth, height=0.44\linewidth]
  \addplot[only marks, mark=*, mark size=1.3pt, omDef, fill opacity=0.6] table[x=gdp, y=labsh, col sep=comma] {figdata/bachelor-1/16-factor-markets/out/labsh_2019.csv};
  \addplot[only marks, mark=none, nodes near coords, point meta=explicit symbolic, nodes near coords style={font=\scriptsize, anchor=west}]
    table[x=gdp, y=labsh, meta=country, col sep=comma] {figdata/bachelor-1/16-factor-markets/out/labsh_labels.csv};
\end{axis}
\end{tikzpicture}
\end{openingchart}

The median country pays $0.521$ of its GDP to labour; a tenth pay less than
$0.329$ and a tenth more than $0.642$. % figdata: labsh_stats
The share is barely related to income: the correlation with the logarithm
of income per person is $0.12$. % figdata: labsh_stats
The lowest shares are in oil exporters such as Qatar, at $0.177$, where most
income is a return to oil in the ground. % figdata: labsh_points
The chart's shares are measured with care: a naive share, counting only
the wages of employees, looks very different, as
\cref{sec:b1:factor-markets:rent} shows.

\begin{recall}
Labour demand from the \omterm{def:b1:factor-markets:vmp}{value of the marginal product}, and the minimum wage
in a competitive market and with one large employer, were met at school.
Factor demands of a price-taking firm are in \cref{ch:b1:profit-maximisation},
\omterm{def:b1:competitive-industry:rent}{economic rent} in \cref{ch:b1:competitive-industry}, and monopoly in
\cref{ch:b1:monopoly-pricing}.
\end{recall}

\section{Derived demand}\label{sec:b1:factor-markets:derived}

\begin{definition}[Derived demand]\label{def:b1:factor-markets:derived-demand}
\index{derived demand}
The demand for a factor of production is a \emph{derived demand}: firms want
the factor not for itself but for the output it helps produce, so its
demand depends on the demand for that output.
\end{definition}

\begin{definition}[Marginal revenue product]\label{def:b1:factor-markets:mrp}
\index{marginal revenue product}
The \emph{marginal revenue product} of a factor is the extra revenue from
one more unit of it: $\mathrm{MRP}_L = MR(q)\cdot\partial f/\partial L$, the marginal revenue of output
times the factor's \omterm{def:b1:technology:marginal-product}{marginal product}.
\end{definition}

\begin{definition}[Value of the marginal product]\label{def:b1:factor-markets:vmp}
\index{value of the marginal product}
For a firm that takes the output price $p$ as given, $MR = p$, and the
\omterm{def:b1:factor-markets:mrp}{marginal revenue product} is the \emph{value of the \omterm{def:b1:technology:marginal-product}{marginal product}},
$\mathrm{VMP}_L = p\,\partial f/\partial L$.
\end{definition}

A firm that takes the wage as given hires until $\mathrm{MRP}_L = w$; its labour
demand is the downward-sloping part of the MRP curve. How elastic that
demand is follows a set of rules often given the label \emph{Marshall's
rules}: the demand for a factor is more elastic when the demand for the
output is more elastic, when other factors substitute for it more easily,
when their supply is more elastic, and (usually) when the factor's share of
cost is larger. For a firm with constant returns and an \omterm{def:b1:technology:sigma}{elasticity of
substitution} $\sigma$, the long-run elasticity of labour demand is
\[
\eta_L = -\bigl[(1 - s_L)\sigma + s_L|\varepsilon|\bigr],
\]
with $s_L$ labour's share of cost and $\varepsilon$ the elasticity of output demand: a
\omterm{def:b1:income-and-substitution-effects:effects}{substitution effect} at given output plus a scale effect. We prove the first
part.

\begin{proposition}[The substitution effect]\label{prop:b1:factor-markets:substitution}
For the Cobb--Douglas technology $q = L^{\alpha}K^{1-\alpha}$, the cost-minimising labour
demand at a given output has elasticity $-(1 - \alpha) = -(1 - s_L)$ with respect
to the wage.
\end{proposition}

\begin{proof}
The tangency gives $K/L = \frac{1 - \alpha}{\alpha}\frac wr$, so $q = L\,(K/L)^{1-\alpha}$ and
$L^* = q\left(\frac{\alpha}{1 - \alpha}\frac rw\right)^{1-\alpha}$, whose elasticity with respect to $w$ is $-(1 - \alpha)$.
Labour's share of cost is $\alpha$ (\cref{ch:b1:cost-minimisation}), and $\sigma = 1$.
\end{proof}

\section{One buyer}\label{sec:b1:factor-markets:monopsony}

A firm that is the only employer in a valley faces the whole upward-sloping
labour supply $w(L)$: to hire more workers it must raise the wage, for the
new workers and, with one wage for all, for those already employed.

\begin{definition}[Monopsony]\label{def:b1:factor-markets:monopsony}
\index{monopsony}\index{single buyer}
A \emph{monopsony} is a market with a single buyer; in a labour market, a
single employer who faces an upward-sloping labour supply and sets the
wage.
\end{definition}

\begin{definition}[Marginal cost of labour]\label{def:b1:factor-markets:mcl}
\index{marginal cost of labour}
The \emph{marginal cost of labour} is the increase in the wage bill from one
more worker: $\mathrm{MCL} = \frac{\dd}{\dd L}\bigl(w(L)L\bigr) = w(L) + Lw'(L)$, above the wage when supply
slopes up.
\end{definition}

The monopsonist hires where $\mathrm{MRP} = \mathrm{MCL}$ and pays the wage on the supply
curve at that employment, below the \omterm{def:b1:factor-markets:mrp}{marginal revenue product}.

\begin{definition}[Markdown]\label{def:b1:factor-markets:markdown}
\index{markdown}
The \emph{markdown} of a wage is the gap between the \omterm{def:b1:factor-markets:mrp}{marginal revenue
product} and the wage, as a share of the wage: $(\mathrm{MRP} - w)/w$.
\end{definition}

\begin{proposition}[The markdown rule]\label{prop:b1:factor-markets:markdown}
At the monopsonist's choice, $\dfrac{\mathrm{MRP} - w}{w} = \dfrac1{\varepsilon_L}$, where
$\varepsilon_L = \frac{\dd L}{\dd w}\frac wL$ is the elasticity of labour supply to the firm.
\end{proposition}

\begin{proof}
$\mathrm{MCL} = w + Lw'(L) = w\bigl(1 + \frac Lw\frac{\dd w}{\dd L}\bigr) = w(1 + 1/\varepsilon_L)$. Setting it equal to
MRP gives $\mathrm{MRP}/w - 1 = 1/\varepsilon_L$.
\end{proof}

The rule is the mirror of the inverse-elasticity rule of
\cref{prop:b1:monopoly-pricing:inverse-elasticity}: the less easily workers
can leave, the larger the \omterm{def:b1:factor-markets:markdown}{markdown}. A firm facing an elasticity of 4 pays
$1/1.25 = 80\%$ of the \omterm{def:b1:factor-markets:mrp}{marginal revenue product}.

\begin{omfigure}
\begin{tikzpicture}
\begin{axis}[omrecord, grid=none, xlabel={Workers $L$}, ylabel={Dalens an hour}, xmin=0, xmax=240, ymin=0, ymax=30,
  width=0.56\linewidth, height=0.42\linewidth]
  \fill[omProp!35] (axis cs:100,11.667) -- (axis cs:100,18.333) -- (axis cs:150,15) -- cycle;
  \addplot[demand, domain=0:240] {25 - x/15};
  \addplot[supply, domain=0:240] {5 + x/15};
  \addplot[shifted, domain=0:150] {5 + 2*x/15};
  \addplot[guide] coordinates {(0,11.667) (100,11.667) (100,0)};
  \addplot[guide] coordinates {(100,18.333) (0,18.333)};
  \addplot[guide] coordinates {(150,15) (150,0)};
  \node[font=\footnotesize, anchor=west] at (axis cs:200,13.5) {MRP};
  \node[font=\footnotesize, anchor=south east] at (axis cs:236,21) {supply $w(L)$};
  \node[font=\footnotesize, anchor=south east] at (axis cs:138,25) {MCL};
\end{axis}
\end{tikzpicture}

{\small The mill of the weekend problem: supply $w = 5 + L/15$ (red),
\omterm{def:b1:factor-markets:mrp}{marginal revenue product} $25 - L/15$ (blue), \omterm{def:b1:factor-markets:mcl}{marginal cost of labour}
$5 + 2L/15$ (dashed). The mill hires 100 workers at 11.67 an hour, whose MRP is
18.33; a competitive market would employ 150 at 15. The triangle is the
\omterm{def:b1:partial-equilibrium-welfare:dwl}{deadweight loss}.}
\end{omfigure}

\section{A wage floor under one buyer}\label{sec:b1:factor-markets:floor}

A legal minimum wage $\bar w$ changes the monopsonist's \omterm{def:b1:factor-markets:mcl}{marginal cost of
labour}. Up to the number of workers willing to work at $\bar w$, each extra
worker costs exactly $\bar w$: the firm no longer has to raise everyone's wage
to hire one more.

\begin{proposition}[A floor can raise employment]\label{prop:b1:factor-markets:floor}
Suppose MRP falls and MCL rises with $L$. Let $L_m$, $w_m$ be the \omterm{def:b1:factor-markets:monopsony}{monopsony}
employment and wage and $L_c$, $w_c$ the competitive ones. A floor $\bar w$ with $w_m < \bar w \le w_c$ raises employment to the
labour supplied at $\bar w$, which is largest, $L_c$, at $\bar w = w_c$. A floor above
$w_c$ sets employment on the MRP curve, which falls as the floor rises, and
it lowers employment below $L_m$ once $\bar w$ exceeds $\mathrm{MRP}(L_m)$.
\end{proposition}

\begin{proof}
Under the floor the \omterm{def:b1:factor-markets:mcl}{marginal cost of labour} is $\bar w$ up to $L^s(\bar w)$ and
$\mathrm{MCL}(L)$ beyond. If $\bar w \le w_c$, $\mathrm{MRP} \ge \bar w$ for all $L \le L^s(\bar w)$, so the firm hires
all $L^s(\bar w)$ workers. Beyond that point $L > L^s(\bar w) \ge L_m$, and since
$\mathrm{MCL} - \mathrm{MRP}$ is increasing and zero at $L_m$, $\mathrm{MCL}(L) > \mathrm{MRP}(L)$: it hires no
more. $L^s$ increases with $\bar w$. If $\bar w > w_c$, the firm can hire as many
as it wants at $\bar w$ and hires until $\mathrm{MRP} = \bar w$, which gives an $L$ falling in
$\bar w$, below $L_m$ when $\bar w > \mathrm{MRP}(L_m)$.
\end{proof}

In the mill town a floor of 13 raises employment from 100 to 120, a floor of
15 to 150, and a floor of 20 cuts it to 75; floors between 11.67 and 18.33
all raise it. The result needs the floor in the right range, and a real
labour market has many firms with different \omterm{def:b1:factor-markets:markdown}{markdowns}, so one floor cannot
be right for all. Whether actual minimum wages raised or lowered employment
is the subject of a debate in Year~3.

\begin{model}{Monopsony}
\assumes One employer, or a few, facing an upward-sloping labour supply,
and one wage for all workers in a job.
\predicts A wage below the \omterm{def:b1:factor-markets:mrp}{marginal revenue product} with a \omterm{def:b1:factor-markets:markdown}{markdown} of
$1/\varepsilon_L$; employment below the competitive level; employment that rises with
a wage floor up to the competitive wage.
\tested Firms differ in what they pay similar workers: firm-specific
premiums explain about $20\%$ of the variation of wages, and wages rise
with the firm's value added per worker with an elasticity of 0.05 to 0.15
(Card, Cardoso, Heining and Kline 2018). % ledger: b1:card2018
Across more than 8,000 US local labour markets, going from the 25th to the
75th percentile of employer concentration is associated with posted wages
5 to $17\%$ lower (Azar, Marinescu and Steinbaum 2022). % ledger: b1:azar2022
\breaks With many employers competing for the same workers; when the
employer can pay each worker a different wage, as in first-degree \omterm{def:b1:monopoly-pricing:discrimination}{price
discrimination}, and so hires the competitive number.
\end{model}

\section{Rent, capital and factor shares}\label{sec:b1:factor-markets:rent}

A factor in fixed supply --- land, a licence, a rare talent --- earns an
\omterm{def:b1:competitive-industry:rent}{economic rent} (\cref{def:b1:competitive-industry:rent}): its price is set
by demand alone and its supply does not respond. Capital is different: it
is produced, wears out, and can be rented or bought.

\begin{definition}[Rental price of capital]\label{def:b1:factor-markets:rental-price}
\index{rental price of capital}
The \emph{rental price of capital} is the price per period at which a unit
of capital (a machine) can be hired:
\[
R = P_K(r + \delta) - \dot P_K,
\]
where $P_K$ is the price of the machine, $r$ the interest rate, $\delta$ its rate of
depreciation and $\dot P_K$ the change in its price over the period.
\end{definition}

The formula follows from no \omterm{def:b1:price-theory-at-work:arbitrage}{arbitrage}, as for housing in
\cref{def:b1:price-theory-at-work:rental-yield}: an owner who buys the
machine at $P_K$, rents it out for $R$ and sells what is left at the end
earns $R + (1 - \delta)(P_K + \dot P_K) - P_K \approx R - \delta P_K + \dot P_K$, which must equal the interest $rP_K$
he could have earned (ignoring the small product $\delta\dot P_K$). A machine of
10,000 dalens at $r = 4\%$, $\delta = 10\%$, whose price falls $2\%$ a year, rents for
$10000 \times 0.14 + 200 = 1600$ dalens a year: a falling price is a cost to its owner.
A price-taking firm rents capital until its value of \omterm{def:b1:technology:marginal-product}{marginal product}
equals $R$.

\begin{omfigure}
\begin{tikzpicture}
\begin{axis}[omrecord, grid=none, ybar stacked, bar width=26pt, xmin=0.4, xmax=3.6, ymin=0, ymax=1800, xtick={1,2,3},
  xticklabels={{$r = 4\%$, $\delta = 10\%$, no price change}, {$\dot P_K/P_K = -2\%$}, {$\dot P_K/P_K = +2\%$}},
  xticklabel style={font=\footnotesize, text width=3.2cm, align=center}, ylabel={Dalens a year},
  width=0.7\linewidth, height=0.38\linewidth, legend style={font=\footnotesize, at={(0.5,-0.32)}, anchor=north, legend columns=3}]
  \addplot[fill=omDef!40, draw=omDef] coordinates {(1,400) (2,400) (3,400)};
  \addplot[fill=omThm!35, draw=omThm] coordinates {(1,1000) (2,1000) (3,800)};
  \addplot[fill=omProp!40, draw=omProp] coordinates {(1,0) (2,200) (3,0)};
  \legend{Interest, Depreciation net of any price rise, Price fall}
\end{axis}
\end{tikzpicture}

{\small The rental price of a machine costing 10,000 dalens as the sum of its
costs to the owner: 1,400 a year with a steady price, 1,600 when its price
falls $2\%$ a year, 1,200 when it rises $2\%$.}
\end{omfigure}

\begin{definition}[Factor share]\label{def:b1:factor-markets:factor-share}
\index{factor share}\index{labour share}
A \emph{factor share} is the part of an economy's income paid to a factor;
the \emph{labour share} is total labour compensation divided by GDP (often
measured at basic prices, without taxes on products).
\end{definition}

With price-taking firms and constant returns, each factor is paid its
\omterm{def:b1:technology:marginal-product}{marginal product} and the shares exhaust income; with a Cobb--Douglas
technology the shares are the exponents, whatever the factor prices. With
an \omterm{def:b1:technology:sigma}{elasticity of substitution} below one, as \cref{ch:b1:technology} found, a rising wage
relative to the rental price raises labour's share. \omterm{def:b1:monopoly-pricing:markup}{Markups} and \omterm{def:b1:factor-markets:markdown}{markdowns}
lower both \omterm{def:b1:factor-markets:factor-share}{factor shares} and leave a share of pure profit.

\begin{measure}[Labour's share and the self-employed]
A farmer working her own land earns one income for her labour, her land
and her tools together, and national accounts record it as ``mixed
income''. Counting only employees' compensation as labour income makes
labour's share very low where most people are self-employed: such naive
shares range from about $0.05$ to $0.80$ across countries, while shares that
count part of the self-employed's income as labour lie between $0.65$ and $0.80$
for most countries (Gollin 2002). % ledger: b1:gollin2002
The Penn World Table chooses among four adjustments --- all mixed income
to labour, the same \omterm{def:b1:factor-markets:factor-share}{labour share} for mixed income as for the rest of the
economy, the same wage for the self-employed as for employees, all of
agriculture's value added to labour --- because mixed income is recorded
for only about 60 countries; employees are 85 to $95\%$ of those working in
advanced economies but below $50\%$, and as low as $4\%$, in many emerging ones.
In 2005 the naive share averaged $42\%$ and ranged from $5\%$ (Nigeria) to $89\%$
(Bhutan) (Inklaar and Timmer 2013). % ledger: b1:pwt-labsh
\end{measure}

\begin{omchart}{Labour's share of GDP against the self-employed share of
employment, 134 countries, 2019.}{Penn World Table 11.0 (labsh); World
Bank, World Development Indicators (SL.EMP.SELF.ZS)}
\begin{tikzpicture}
\begin{axis}[omscatter, xmin=0, xmax=100, ymin=0, ymax=0.85, xlabel={Self-employed (\% of employment)},
  ylabel={Labour share of GDP}, width=0.66\linewidth, height=0.4\linewidth]
  \addplot[only marks, mark=*, mark size=1.3pt, omDef, fill opacity=0.6] table[x=self, y=labsh, col sep=comma] {figdata/bachelor-1/16-factor-markets/out/labsh_self.csv};
  \addplot[only marks, mark=none, nodes near coords, point meta=explicit symbolic, nodes near coords style={font=\scriptsize, anchor=west}]
    table[x=self, y=labsh, meta=country, col sep=comma] {figdata/bachelor-1/16-factor-markets/out/labsh_labels.csv};
  % figdata: labsh_points (United States: self 6.5, labsh 0.592)
  \draw[black!60] (axis cs:6.5,0.592) -- (axis cs:14,0.78) node[anchor=west, font=\scriptsize, black] {United States};
\end{axis}
\end{tikzpicture}
\end{omchart}

After the adjustment, labour's share is nearly unrelated to how many
people are self-employed: the correlation across 134 countries is $-0.12$.
% figdata: labsh_stats
India, where $77.2\%$ of workers are self-employed, has a share of $0.533$; the
United States, with $6.5\%$, has $0.592$. % figdata: labsh_points
The adjustment does what it was meant to do; whether it does it right in
each country is harder to know.

\section*{Exercises}\label{sec:b1:factor-markets:exercises}

\begin{exercise}[$\star$]\label{exo:b1:factor-markets:vmp-ex}
A price-taking firm produces $q = 10L^{0.5}$ and sells at $p = 4$. Find its VMP of
labour and the number of workers it hires at a wage of 5.
\end{exercise}

\begin{exercise}[$\star$]\label{exo:b1:factor-markets:read}
From the charts' table, France's \omterm{def:b1:factor-markets:factor-share}{labour share} is $0.642$ with $12.1\%$ of
workers self-employed, and Nigeria's $0.312$ with $86.6\%$. % figdata: labsh_points
What would a naive share, counting only employees, probably show for
Nigeria, and why?
\end{exercise}

\begin{exercise}[$\star$]\label{exo:b1:factor-markets:mono-ex}
Labour supply to a firm is $w = 2 + 0.1L$ and $\mathrm{MRP} = 20 - 0.1L$. Find the
\omterm{def:b1:factor-markets:monopsony}{monopsony} wage and employment, and the competitive ones.
\end{exercise}

\begin{exercise}[$\star$]\label{exo:b1:factor-markets:floor12}
In the market of \cref{exo:b1:factor-markets:mono-ex}, predict employment
after a wage floor of 12.
\end{exercise}

\begin{exercise}[$\star\star$]\label{exo:b1:factor-markets:best-floor}
In the same market, find the floor that maximises employment and the
number of jobs it adds.
\end{exercise}

\begin{exercise}[$\star\star$]\label{exo:b1:factor-markets:markdown4}
The elasticity of labour supply to a firm is 4. Compute the \omterm{def:b1:factor-markets:markdown}{markdown} and
the wage as a share of the MRP.
\end{exercise}

\begin{exercise}[$\star\star$]\label{exo:b1:factor-markets:derived-ex}
A Cobb--Douglas firm with constant returns has labour's share of cost 0.7
and faces an output demand of elasticity $-2$. Compute its long-run
elasticity of labour demand and predict what happens if its output demand
becomes less elastic.
\end{exercise}

\begin{exercise}[$\star\star$]\label{exo:b1:factor-markets:rental}
Compute the rental price of a machine costing 10,000 with $r = 4\%$, $\delta = 10\%$,
when its price falls $2\%$ a year. Why do computers rent for so much of
their price?
\end{exercise}

\begin{exercise}[$\star\star$]\label{exo:b1:factor-markets:corr}
The correlation of the \omterm{def:b1:factor-markets:factor-share}{labour share} with the self-employed share across
134 countries is $-0.12$. % figdata: labsh_stats
What would you expect it to be for naive shares, and why?
\end{exercise}

\begin{exercise}[$\star\star\star$]\label{exo:b1:factor-markets:derive}
Derive the \omterm{def:b1:factor-markets:markdown}{markdown} rule, and show that a monopsonist facing a perfectly
elastic labour supply pays the MRP.
\end{exercise}

\begin{exercise}[$\star\star\star$]\label{exo:b1:factor-markets:above}
In the market of \cref{exo:b1:factor-markets:mono-ex}, a floor of 13 is set.
Compare employment with the \omterm{def:b1:factor-markets:monopsony}{monopsony} and with a competitive market with
the same supply and MRP curves.
\end{exercise}

\begin{exercise}[$\star\star\star$]\label{exo:b1:factor-markets:separate}
A rising labour supply to a whole industry raises wages as the industry
grows even if every firm is a price taker. What evidence would tell a
monopsonist apart from such a competitive industry?
\end{exercise}

\section{Problem: The Mill Town}\label{sec:b1:factor-markets:problem}

\begin{problem}[{Weekend problem --- the employment-maximising floor}]\label{pb:b1:factor-markets:mill}
A valley has one mill (invented). The labour supply from the valley is
$w = 5 + L/15$ dalens an hour for $L$ workers, and the mill's \omterm{def:b1:factor-markets:mrp}{marginal revenue
product} is $25 - L/15$. Sums below are in dalens an hour.

\medskip\noindent\textbf{Part I --- The competitive benchmark.}
\begin{enumerate}
  \item Find the wage and employment if the mill took the wage as given.
  \item Compute the elasticity of labour supply there.
  \item Compute the workers' surplus (the area between the wage and the supply curve).
  \item Compute the mill's surplus (the area between MRP and the wage).
  \item Compute \omterm{def:b1:partial-equilibrium-welfare:surplus}{total surplus}.
\end{enumerate}

\medskip\noindent\textbf{Part II --- The mill sets the wage.}
\begin{enumerate}[resume]
  \item Write the wage bill and the \omterm{def:b1:factor-markets:mcl}{marginal cost of labour}.
  \item Find the \omterm{def:b1:factor-markets:monopsony}{monopsony} employment and wage.
  \item Compute the MRP at that employment and the \omterm{def:b1:factor-markets:markdown}{markdown}.
  \item Check the \omterm{def:b1:factor-markets:markdown}{markdown} against the supply elasticity.
  \item Compute the mill's surplus.
  \item Compute the workers' surplus.
  \item Compute the \omterm{def:b1:partial-equilibrium-welfare:dwl}{deadweight loss}.
\end{enumerate}

\medskip\noindent\textbf{Part III --- A wage floor.}
\begin{enumerate}[resume]
  \item Find employment under a floor of 13.
  \item Under a floor of 15.
  \item Under a floor of 17.
  \item Find the range of floors that raise employment above the \omterm{def:b1:factor-markets:monopsony}{monopsony} level.
  \item Find employment under a floor of 20.
  \item Compute the workers' and the mill's surplus under a floor of 15, and compare with Part II.
  \item Explain why the floor of 15 raises employment here, while a floor of 15 would lower it in a competitive market whose wage is 12.
\end{enumerate}

\medskip\noindent\textbf{Part IV --- A second mill.} An identical mill opens;
each has the \omterm{def:b1:factor-markets:mrp}{marginal revenue product} $25 - L_i/15$ and both hire from the
same valley, choosing their employment and taking the other's as given.
\begin{enumerate}[resume]
  \item Write mill 1's \omterm{def:b1:factor-markets:mcl}{marginal cost of labour}.
  \item Find the symmetric \omterm{def:b1:models-and-data:equilibrium}{equilibrium}: employment at each mill, in all, and the wage.
  \item Compute the MRP and the \omterm{def:b1:factor-markets:markdown}{markdown}.
  \item What elasticity of labour supply does each mill face? Check the \omterm{def:b1:factor-markets:markdown}{markdown} rule.
  \item Compare with a competitive market with two mills.
  \item State the result: the employment-maximising wage floor in the one-mill town, and the workers it adds.
\end{enumerate}
\end{problem}
